\documentclass{Phyport}

\usepackage{hyperref}
\hypersetup{
    colorlinks=true,
}
% 选用宏（如有额外需要的宏包请在此添加）
\usepackage{float}       % 提供 [H] 选项
\usepackage{booktabs}    % 三线表
\usepackage{amsmath}     % 数学公式
\usepackage{caption}     % 标题设置
\usepackage{siunitx}     % SI 单位排版
 %选用宏

\exname{\small 铁磁材料的磁滞回线和基本磁化曲线} %实验名称
\extable{9} %实验桌号
\instructor{陈飞燕} %指导教师
\class{电科2502} %班级
\name{管云哲} %姓名
\stuid{3250104519} %学号

\nyear{2026} %年
\nmonth{10} %月
\nday{20} %日
\nweekday{二} %星期几

\redate{} %如有实验补做，补做日期
\resitu{} %情况说明：

\begin{document}
\setcounter{page}{1}
\makecover

\section{预习报告（10分）}
（注：将已经写好的“物理实验预习报告”内容拷贝过来）

\subsection{实验综述（5分）}
（自述实验现象、实验原理和实验方法，包括必要的光路图、电路图、公式等。不超过500字。）

铁磁材料内部存在自发磁化的小区域——磁畴。外磁场增强时，磁畴转动与畴壁迁移使 $B$ 随 $H$ 非线性增大，$H$ 足够大后达到磁饱和，对应饱和磁感应强度 $B_s$。减小 $H$ 时 $B$ 不沿原路返回：$H=0$ 时保留剩磁 $B_r$，需反向磁场 $H_c$（矫顽力）才能使 $B=0$。$H$ 循环变化时 $B$ 形成闭合磁滞回线；不同幅值回线顶点的连线为基本磁化曲线，其斜率给出磁导率 $\mu=B/H$。动态回线还包含涡流损耗，面积大于静态回线并随频率变化；磁滞损耗正比于回线面积。

\begin{figure}[H]
    \centering
    \includegraphics[width=0.34\linewidth]{bh_curve.jpg}
    \caption{起始磁化曲线与磁滞回线}
    \label{fig:bh-curve}
\end{figure}

实验采用转换法。环形样品上绕励磁线圈 $N_1$ 和测量线圈 $N_2$；$N_1$ 中通交流电流 $I_1$，由安培环路定理 $H=N_1I_1/L$。取样电阻 $R_1$ 上的电压 $U_x=I_1R_1$ 接入示波器 X 轴，故 $U_x\propto H$。$N_2$ 的感应电动势 $E_2=N_2S\,\mathrm{d}B/\mathrm{d}t$ 经 $R_2C$ 积分电路（$R_2\gg 1/2\pi fC$）积分，电容电压 $U_y=N_2SB/(R_2C)$ 接入 Y 轴，故 $U_y\propto B$。示波器置 X-Y 模式即可显示磁滞回线，由
\begin{equation*}
    H=\frac{N_1U_x}{LR_1},\qquad B=\frac{R_2CU_y}{N_2S},
\end{equation*}
或按灵敏度 $S_x$、$S_y$ 与格数 $X$、$Y$ 定标求 $H$、$B$。

\begin{figure}[H]
    \centering
    \includegraphics[width=0.72\linewidth]{bh_circuit.jpg}
    \caption{示波器显示 $B$-$H$ 曲线的原理线路}
    \label{fig:bh-circuit}
\end{figure}

实验先升压至饱和调好图形，再单调降压至零退磁；以 0.5\,V 步长升压，记录各回线顶点得基本磁化曲线；保持 3\,V 逐点读取饱和回线坐标，求 $B_s$、$B_r$、$H_c$ 及回线面积，并比较不同频率与样品。

\subsection{实验重点（3分）}
（简述本实验的学习重点，不超过100字。）

\begin{enumerate}
    \item 理解磁滞、剩磁、矫顽力及基本磁化曲线的物理意义；
    \item 掌握示波器 X-Y 法（转换法）显示动态磁滞回线的原理；
    \item 学会退磁、定标及由回线求 $B_s$、$B_r$、$H_c$ 与磁滞损耗。
\end{enumerate}

\subsection{实验难点（2分）}
（简述本实验的实现难点，不超过100字。）

\begin{enumerate}
    \item 退磁须单调降压至零，否则剩磁使曲线偏离原点；
    \item X-Y 增益、耦合与光斑位置需协调，否则回线畸变或显示不全；
    \item 光标读取交点与顶点误差较大，直接影响 $B_r$、$H_c$ 和面积；
    \item 积分条件 $R_2\gg 1/(2\pi fC)$ 不满足时回线出现畸变。
\end{enumerate}

\section{实验操作与原始数据（50分）}
（将有老师签名的“自备数据记录草稿纸”的扫描或手机拍摄图粘贴在下方，完整保留姓名，学号，教师签字和日期。按老师的要求拍摄实验过程图，附在此处）

实验按讲义步骤进行：先将信号源输出调至最小，接通示波器与磁特性测量仪并把光点调到屏幕中心；对样品单调升压至磁化饱和，调节示波器 X、Y 增益使回线大小合适并锁定；再单调降压至零，使样品充分退磁。随后从 $U=0$ 开始逐级升压，记录各磁滞回线顶点的坐标，得到基本磁化曲线；保持励磁电压 $U=3.00$\,V，沿饱和回线逐点读取上、下支路的 $U_x$、$U_y$ 坐标。原始记录如下。

\begin{figure}[H]
    \centering
    \includegraphics[width=0.58\linewidth]{实验数据.jpg}
    \caption{原始数据记录（姓名、学号及教师签字）}
    \label{fig:raw-data}
\end{figure}

\section{结果与分析（30分）}
\subsection{数据处理与结果（20分）}
（列出数据表格、选择适合的数据处理方法、写出测量或计算结果。）

实验使用样品 1（环形铁氧体），其参数为
\begin{equation*}
    N_1=N_2=150,\qquad L=0.130\ \mathrm{m},\qquad S=1.24\times10^{-4}\ \mathrm{m^2},
\end{equation*}
电路参数为
\begin{equation*}
    R_1=2.5\ \Omega,\qquad R_2=1.0\times10^{4}\ \Omega,\qquad C=3.0\times10^{-6}\ \mathrm{F},\qquad f=130\ \mathrm{Hz}.
\end{equation*}

由取样电阻与积分电路关系
\begin{equation}
    H=\frac{N_1U_x}{R_1L}=461.5\,U_x,\qquad
    B=\frac{R_2CU_y}{N_2S}=1.613\,U_y
    \label{eq:convert}
\end{equation}
（式中 $U_x$、$U_y$ 以 V 为单位，$H$ 的单位为 A/m，$B$ 的单位为 T），并取
\begin{equation*}
    \mu=\frac{B}{H},\qquad \mu_r=\frac{\mu}{\mu_0},\qquad \mu_0=4\pi\times10^{-7}\ \mathrm{H/m}.
\end{equation*}

\subsubsection*{1. 基本磁化曲线与 $\mu$-$H$ 关系}

将各励磁电压下回线顶点的读数代入式~\eqref{eq:convert}，结果如表~\ref{tab:basic} 所示。

\begin{table}[H]
    \centering
    \caption{基本磁化曲线与 $\mu$-$H$ 曲线数据}
    \label{tab:basic}
    \begin{tabular}{ccccccc}
        \toprule
        $U$/V & $U_x$/mV & $U_y$/mV & $H$/(A$\cdot$m$^{-1}$) & $B$/mT & $\mu$/($10^{-3}$\,H$\cdot$m$^{-1}$) & $\mu_r$ \\
        \midrule
        0.5 & 47.58  & 39.35  & 21.96  & 63.47  & 2.890 & 2300 \\
        1.0 & 95.52  & 79.64  & 44.09  & 128.45 & 2.914 & 2319 \\
        1.5 & 143.10 & 120.20 & 66.05  & 193.87 & 2.935 & 2336 \\
        2.0 & 192.60 & 159.20 & 88.89  & 256.77 & 2.889 & 2299 \\
        2.5 & 284.10 & 199.80 & 131.12 & 322.26 & 2.458 & 1956 \\
        3.0 & 702.60 & 237.40 & 324.28 & 382.90 & 1.181 & 940  \\
        \bottomrule
    \end{tabular}
\end{table}

\begin{figure}[H]
    \centering
    \includegraphics[width=0.78\linewidth]{basic_curve.pdf}
    \caption{样品 1 的基本磁化曲线与 $\mu_r$-$H$ 关系}
    \label{fig:basic}
\end{figure}

由表~\ref{tab:basic} 和图~\ref{fig:basic} 可见：$H$ 较小时 $B$ 随 $H$ 近似线性增大，$\mu_r$ 基本保持在 $2.3\times10^{3}$ 左右；$H$ 超过约 $130$\,A/m 后 $B$ 增长变缓并趋于饱和，$\mu_r$ 迅速下降，在 $H=324.28$\,A/m 时约为 $940$。这说明铁磁材料的磁导率不是常量，而是随磁化状态变化。

\subsubsection*{2. 饱和磁滞回线与特征参量}

$U=3.00$\,V 时沿回线上、下支路逐点读取的 48 组数据经式~\eqref{eq:convert} 换算后列于表~\ref{tab:loop}，据此绘出饱和磁滞回线（图~\ref{fig:loop}）。

\begin{table}[H]
    \centering
    \caption{饱和磁滞回线数据（$U=3.00$\,V）}
    \label{tab:loop}
    \footnotesize
    \setlength{\tabcolsep}{4pt}
    \begin{tabular}{c r r r r c r r r r}
        \toprule
        No & $U_x$/mV & $U_y$/mV & $H$/(A$\cdot$m$^{-1}$) & $B$/mT &
        No & $U_x$/mV & $U_y$/mV & $H$/(A$\cdot$m$^{-1}$) & $B$/mT \\
        \midrule
        1  & 702.6  & 237.40  & 324.28  & 382.90  & 25 & 50.0   & $-142.10$ & 23.08   & $-229.19$ \\
        2  & 650.0  & 234.20  & 300.00  & 377.74  & 26 & 0.0    & 162.50    & 0.00    & 262.10   \\
        3  & 600.0  & 231.10  & 276.92  & 372.74  & 27 & 0.0    & $-159.30$ & 0.00    & $-256.94$ \\
        4  & 550.0  & 228.00  & 253.85  & 367.74  & 28 & $-50.0$ & 143.70   & $-23.08$ & 231.77  \\
        5  & 500.0  & 225.40  & 230.77  & 363.55  & 29 & $-50.0$ & $-173.40$ & $-23.08$ & $-279.68$ \\
        6  & 500.0  & 218.60  & 230.77  & 352.58  & 30 & $-100.0$ & 123.40  & $-46.15$ & 199.03  \\
        7  & 450.0  & 223.30  & 207.69  & 360.16  & 31 & $-100.0$ & $-184.30$ & $-46.15$ & $-297.26$ \\
        8  & 450.0  & 210.80  & 207.69  & 340.00  & 32 & $-150.0$ & 95.31   & $-69.23$ & 153.73  \\
        9  & 400.0  & 220.20  & 184.62  & 355.16  & 33 & $-150.0$ & $-195.50$ & $-69.23$ & $-315.32$ \\
        10 & 400.0  & 196.70  & 184.62  & 317.26  & 34 & $-200.0$ & 54.68   & $-92.31$ & 88.19   \\
        11 & 350.0  & 217.00  & 161.54  & 350.00  & 35 & $-200.0$ & $-203.40$ & $-92.31$ & $-328.06$ \\
        12 & 350.0  & 174.90  & 161.54  & 282.10  & 36 & $-246.8$ & 0.00    & $-113.91$ & 0.00  \\
        13 & 300.0  & 210.80  & 138.46  & 340.00  & 37 & $-246.8$ & $-209.60$ & $-113.91$ & $-338.06$ \\
        14 & 300.0  & 115.50  & 138.46  & 186.29  & 38 & $-300.0$ & $-106.20$ & $-138.46$ & $-171.29$ \\
        15 & 250.0  & 204.50  & 115.38  & 329.84  & 39 & $-300.0$ & $-214.30$ & $-138.46$ & $-345.65$ \\
        16 & 250.0  & 6.16    & 115.38  & 9.94    & 40 & $-350.0$ & $-171.80$ & $-161.54$ & $-277.10$ \\
        17 & 246.8  & 0.00    & 113.91  & 0.00    & 41 & $-350.0$ & $-219.00$ & $-161.54$ & $-353.23$ \\
        18 & 200.0  & 200.00  & 92.31   & 322.58  & 42 & $-400.0$ & $-196.80$ & $-184.62$ & $-317.42$ \\
        19 & 200.0  & $-53.12$ & 92.31  & $-85.68$ & 43 & $-400.0$ & $-222.10$ & $-184.62$ & $-358.23$ \\
        20 & 150.0  & 192.10  & 69.23   & 309.84  & 44 & $-450.0$ & $-210.90$ & $-207.69$ & $-340.16$ \\
        21 & 150.0  & $-92.18$ & 69.23  & $-148.68$ & 45 & $-450.0$ & $-226.80$ & $-207.69$ & $-365.81$ \\
        22 & 100.0  & 182.80  & 46.15   & 294.84  & 46 & $-500.0$ & $-218.70$ & $-230.77$ & $-352.74$ \\
        23 & 100.0  & $-123.40$ & 46.15 & $-199.03$ & 47 & $-500.0$ & $-229.90$ & $-230.77$ & $-370.81$ \\
        24 & 50.0   & 173.40  & 23.08   & 279.68  & 48 & $-684.3$ & $-235.90$ & $-315.83$ & $-380.48$ \\
        \bottomrule
    \end{tabular}
\end{table}

\begin{figure}[H]
    \centering
    \includegraphics[width=0.52\linewidth]{hysteresis_loop.pdf}
    \caption{样品 1 的饱和磁滞回线（$U=3.00$\,V）}
    \label{fig:loop}
\end{figure}

由图~\ref{fig:loop} 可直接读出：正向饱和点 $B_{m+}=382.90$\,mT、反向饱和点 $B_{m-}=-380.48$\,mT，取平均得
\begin{equation*}
    B_m=\frac{|B_{m+}|+|B_{m-}|}{2}=381.69\ \mathrm{mT}.
\end{equation*}
$H=0$ 时回线与 $B$ 轴的交点为 $B_{r+}=262.10$\,mT、$B_{r-}=-256.94$\,mT，故
\begin{equation*}
    B_r=\frac{|B_{r+}|+|B_{r-}|}{2}=259.52\ \mathrm{mT}.
\end{equation*}
$B=0$ 时回线与 $H$ 轴的交点为 $H_{c+}=113.91$\,A/m、$H_{c-}=-113.91$\,A/m，故
\begin{equation*}
    H_c=\frac{|H_{c+}|+|H_{c-}|}{2}=113.91\ \mathrm{A/m}.
\end{equation*}

回线所围面积表示单位体积样品在一个磁化周期内消耗的能量（动态测量时还包含涡流损耗）。将回线离散点按上、下支路顺序闭合，用多边形面积公式
\begin{equation*}
    |BH|=\frac{1}{2}\left|\sum_i\left(H_iB_{i+1}-H_{i+1}B_i\right)\right|
\end{equation*}
计算，得
\begin{equation*}
    |BH|\approx 1.29\times10^{2}\ \mathrm{J/m^3}.
\end{equation*}

\subsection{误差分析（10分）}
（运用测量误差、相对误差或不确定度等分析实验结果，写出完整的结果表达式，并分析误差原因。）

由式~\eqref{eq:convert}，$N_1$、$N_2$ 为整数且视为准确值，$H$、$B$ 的相对不确定度分别为
\begin{equation*}
    \frac{u_H}{H}=\sqrt{\left(\frac{u_{U_x}}{U_x}\right)^2+\left(\frac{u_{R_1}}{R_1}\right)^2+\left(\frac{u_L}{L}\right)^2},\qquad
    \frac{u_B}{B}=\sqrt{\left(\frac{u_{U_y}}{U_y}\right)^2+\left(\frac{u_{R_2}}{R_2}\right)^2+\left(\frac{u_C}{C}\right)^2+\left(\frac{u_S}{S}\right)^2}.
\end{equation*}
取示波器光标读数误差 $u_{U_x}=u_{U_y}=0.5$\,mV，电阻误差 $1\%$，电容误差 $5\%$，磁路长度误差 $0.5$\,mm，截面积误差 $1\%$，代入特征点数据得
\begin{align*}
    H_c &= (1.14\pm0.01)\times10^{2}\ \mathrm{A/m}\quad(1.1\%),\\
    B_r &= (2.60\pm0.14)\times10^{-1}\ \mathrm{T}\quad(5.2\%),\\
    B_m &= (3.82\pm0.20)\times10^{-1}\ \mathrm{T}\quad(5.2\%),\\
    |BH| &= (1.29\pm0.07)\times10^{2}\ \mathrm{J/m^3}\quad(5.3\%).
\end{align*}

误差主要来源：
\begin{enumerate}
    \item 示波器光标读数误差。回线有一定线宽，读取顶点及与坐标轴交点时会引入人为判断误差，对 $H_c$、$B_r$ 影响较直接。
    \item 电路元件参数误差。$B$ 正比于 $R_2C$，电容的容差（按 $5\%$ 估计）是 $B$ 的主要误差来源；$H$ 正比于 $1/R_1$，取样电阻的容差也有影响。
    \item 样品参数误差。$L$、$S$ 按平均磁路长度和几何尺寸计算，实际磁化并不均匀，会带来系统误差。
    \item 积分近似误差。本实验 $1/(2\pi fC)=408\ \Omega$，$R_2=10\ \mathrm{k}\Omega\gg 408\ \Omega$，积分条件满足；但积分电路并非理想，仍会带来幅值与相位误差。
    \item 动态测量误差。交流磁化时样品中除磁滞损耗外还有涡流损耗，所得回线面积大于静态回线，且与频率有关。
    \item 数据处理误差。回线面积用有限个离散点作多边形近似，取点间距和支路连接方式都会影响面积结果。
\end{enumerate}

\section{实验探讨与思考题（10分）}
（对实验内容、现象和过程的小结，不超过100字。并解答教材或讲义或老师布置的思考题，请先写题干，再作答。）

实验用示波器 X-Y 法测得铁氧体的基本磁化曲线和饱和磁滞回线。$\mu_r$ 在低场约为 $2.3\times10^{3}$，饱和后降至 $940$；$B_r\approx0.26$\,T，$H_c\approx114$\,A/m，回线面积约 $129$\,J/m$^3$。动态回线包含涡流损耗，故大于静态回线。

\subsection*{思考题 1}
\textbf{题目：}$B\cdot H$（磁能积）的单位是否与能量有关？请给出量纲的推导过程。

\textbf{答：}有关。磁能积的单位是能量密度而非能量本身。由 $B$ 的单位 T $=\mathrm{N/(A\cdot m)}$、$H$ 的单位 A/m，得
\begin{equation*}
    [B\cdot H]=\mathrm{T\cdot A\cdot m^{-1}}
    =\frac{\mathrm{N}}{\mathrm{A\cdot m}}\cdot\frac{\mathrm{A}}{\mathrm{m}}
    =\frac{\mathrm{N}}{\mathrm{m^2}}
    =\frac{\mathrm{N\cdot m}}{\mathrm{m^3}}
    =\frac{\mathrm{J}}{\mathrm{m^3}}.
\end{equation*}
因此 $B\cdot H$ 的量纲为 $\mathrm{J/m^3}$，表示单位体积磁性材料在一个磁化周期中因磁滞效应损耗的能量。

\subsection*{思考题 2}
\textbf{题目：}查阅有关资料，简要回答积分电路的原理与有关应用。

\textbf{答：}积分电路通常由电阻 $R$ 和电容 $C$ 串联组成，输出电压取自电容两端。当输入信号为 $E(t)$ 时，回路电流 $I\approx E/R$（要求 $R\gg1/(\omega C)$），电容电压
\begin{equation*}
    U_{\mathrm{out}}=U_C=\frac{1}{C}\int I\,\mathrm{d}t\approx\frac{1}{RC}\int E(t)\,\mathrm{d}t,
\end{equation*}
即输出电压近似与输入电压对时间的积分成正比。

本实验中，测量线圈的感应电动势 $E_2=N_2S\,\mathrm{d}B/\mathrm{d}t$ 经 $R_2C$ 积分后，输出电压 $U_y\propto B$，从而把磁感应强度的变化率转换为与 $B$ 成正比的电压，使示波器能够直接显示磁滞回线。积分电路还广泛用于波形变换（方波积分为三角波）、模拟计算、峰值检测、传感器信号调理、电荷放大和低通滤波等场合。

\subsection*{思考题 3}
\textbf{题目：}磁性材料内部的涡流损耗的能量是否可以被利用起来？请列出应用实例。

\textbf{答：}可以。涡流损耗本质上是交变磁场在导体内部感应出的涡旋电流在材料电阻上转化为焦耳热的过程。在变压器、电机铁芯中它是有害的，通常用硅钢片叠片、提高材料电阻率或采用铁氧体等方式抑制；但它也可以被利用，例如：
\begin{enumerate}
    \item 感应加热：电磁炉、高频感应熔炼、金属淬火与焊接，利用涡流直接在工件内部产生热量；
    \item 电磁制动：高速列车、过山车和健身器材的涡流刹车，利用涡流产生的阻尼力实现无接触制动；
    \item 涡流阻尼：电表、天平、检流计中的电磁阻尼器，使指针快速稳定；
    \item 无损检测与传感：涡流探伤检测金属表面裂纹，涡流测厚、测距、金属探测器和电磁悬浮等。
\end{enumerate}

{\fangsong 注意事项：
\begin{enumerate}[label=\arabic*.]
    \item 用PDF格式上传“实验报告”，文件名：学生姓名+学号+实验名称+周次。
    \item “实验报告”必须递交在“学在浙大”本课程内对应实验项目的“作业”模块内。
    \item “实验报告”成绩必须在“浙江大学物理实验教学中心网站”-“选课系统”内查询。
    \item 教学评价必须在“浙江大学物理实验教学中心网站”-“选课系统”内进行，学生必须进行教学评价，才能看到实验报告成绩，教学评价须在本次实验结束后3天内进行。
\end{enumerate}}
\begin{center}
    \fbfangsong{浙江大学物理实验教学中心制}
\end{center}
\end{document}
